\section{The \pla{} calculus}\label{sec:calculus}

The repository's semantics document states every definition below next
to an executable block, and the test suite extracts and runs all of
them, so the paper's formal claims cannot drift from the engine.
Figure~\ref{fig:overview} shows the pieces this section and the next
two define: the confidence engine, the symbolic gate, and the learner
that exports its fit back into the same rule vocabulary.

\begin{figure}[t]
\centering
{\linespread{1}\selectfont\small
% System overview (TikZ fragment, \input from a figure environment).
\begin{tikzpicture}[
  font=\small,
  box/.style={draw, rounded corners=2pt, align=center, inner sep=5pt},
  stage/.style={box, fill=black!4, text width=0.26\linewidth},
  io/.style={box, text width=0.20\linewidth},
  side/.style={box, densely dashed, text width=0.26\linewidth},
  arr/.style={-{Stealth[length=2.2mm]}, thick},
  lbl/.style={font=\footnotesize, midway, fill=white, inner sep=1.5pt},
  node distance=5mm and 9mm]

\node[io] (scenario) {\textbf{Scenario (JSON)}\\[1pt]
  facts \\ rules $(\vec{A}{\to}B,\,p,\,w)$ \\
  context sets \\ queries};

\node[stage, right=of scenario] (engine) {\textbf{Confidence engine}\\[2pt]
  antecedents: G\"odel t-norm (min)\\
  context adjust: \texttt{legacy} $|$ \texttt{logit}\\
  aggregate: noisy-OR (default)\\
  Jacobi fixpoint (Prop.~\ref{prop:convergence})};

\node[io, right=of engine] (out) {\textbf{Outputs}\\[1pt]
  confidence $c(B)$ \\ per-rule trace \\ deletion counterfactual};

\node[side, below=of engine] (symbolic) {\textbf{Symbolic layer}\\
  definite clauses, forward chaining\\
  gate: hard $|$ soft};

\node[side, above=of engine] (learner) {\textbf{Learner (Sec.~\ref{sec:learning})}\\
  logistic reduction: $\theta_r,\ w_r(v),\ b$};

\draw[arr] (scenario) -- (engine);
\draw[arr] (engine) -- (out);
\draw[arr] (symbolic) -- (engine) node[lbl, right=2pt] {entailment gate};
\draw[arr] ([xshift=-9mm]learner.south) -- ([xshift=-9mm]engine.north)
  node[lbl, left=2pt] {export weights + prior};
\draw[arr] ([xshift=9mm]engine.north) -- ([xshift=9mm]learner.south)
  node[lbl, right=2pt] {rule firings};
\end{tikzpicture}
}
\caption{\pla{} in one picture. A JSON scenario supplies facts, rules
with base confidences and context weights, and named context sets; the
engine folds fired rules by the chosen aggregator inside a fixpoint
iteration and emits per-rule traces; a crisp definite-clause layer can
gate conclusions; the learner of Section~\ref{sec:learning} fits rule,
context, and intercept weights and exports them back as engine rules.}
\label{fig:overview}
\end{figure}

\subsection{Values, rules, and operators}

Each proposition $x$ carries a confidence $c(x)\in[0,1]$; base facts
have $c=1$. A rule is a triple
$(A_1\wedge\dots\wedge A_n \rightarrow B,\; p,\; w)$ with base
confidence $p\in[0,1]$ and context weights $w:\mathit{Var}\to\mathbb{R}$.
Antecedents combine with the G\"odel t-norm, and a firing rule supports
its head with
\begin{equation}
\text{candidate} \;=\; p_{\text{adj}} \cdot \min_i\, c(A_i).
\label{eq:candidate}
\end{equation}
Candidates for the same head fold through a per-knowledge-base
aggregation operator applied pairwise to the accumulated support $e$ and
the new candidate $n$:
\begin{equation}
\begin{aligned}
\texttt{max}&: \max(e,n) &
\texttt{noisy\_or}&: 1-(1-e)(1-n)\\
\texttt{sum\_cap}&: \min(1, e+n) &
\texttt{logit\_pool}&: \sigma(\logit e + \logit n),
\end{aligned}
\label{eq:aggregators}
\end{equation}
where $\logit p = \ln\frac{p}{1-p}$ and $\sigma(t) = (1+e^{-t})^{-1}$.
The default \texttt{noisy\_or} is identically MYCIN's
parallel-combination rule for positive evidence.

\subsection{Context adjustment}\label{sec:context}

Context adjustment has two modes. The \texttt{legacy} mode multiplies
$p$ by the weight of each active context variable and caps at $1$:
\begin{equation}
p_{\text{adj}} \;=\; \min\Bigl(1,\; p \cdot \prod_{v\ \text{active}} w(v)\Bigr).
\label{eq:legacy}
\end{equation}
The cap destroys information: with $p=0.7$ and weights $1.2$ and $1.5$,
the product $1.26$ caps to $1.0$ --- and so does any stronger
combination, making distinct evidence states indistinguishable. The
\texttt{logit} mode instead treats weights as additive log-odds
deltas:
\begin{equation}
p_{\text{adj}} \;=\; \sigma\Bigl(\logit(p) + \sum_{v\ \text{active}} w(v)\Bigr),
\label{eq:logit}
\end{equation}
which never saturates below $1$, treats strengthening and weakening
(negative weights) symmetrically, keeps $p\in\{0,1\}$ as fixed points,
and matches the naive-Bayes/likelihood-ratio update that is the
probabilistically coherent special case of certainty-factor
reasoning~\cite{heckerman1986}.

\subsection{Inference and convergence}

Inference iterates a Jacobi-style round operator $G$: each round
recomputes every derived confidence from the base facts and the previous
round's values, stopping when no confidence changes by more than
$10^{-9}$.

\begin{proposition}[Convergence]\label{prop:convergence}
Order confidence vectors $x\in[0,1]^V$ pointwise, and let the
aggregation operator be one of \texttt{max}, \texttt{noisy\_or},
\texttt{sum\_cap}. Each of these is monotone in both arguments and the
candidate map \eqref{eq:candidate} is monotone in $x$, so $G$ is a
monotone map on the complete lattice $[0,1]^V$. The start vector $x_0$
(base facts) satisfies $x_0 \le G(x_0)$, hence the trajectory
$x_0 \le G(x_0) \le G^2(x_0) \le \dots$ is pointwise non-decreasing and
bounded by $1$; each coordinate converges, the round-max delta tends to
zero, and by Kleene's fixpoint theorem the limit is the least fixpoint
of $G$ above $x_0$.
\end{proposition}

\texttt{logit\_pool} is deliberately excluded from
Proposition~\ref{prop:convergence}: log-odds pooling treats a candidate
below $1/2$ as \emph{negative} evidence that lowers the accumulated
support, so the operator is not monotone, and its zero-support
convention (a first candidate replaces the empty accumulator) is
discontinuous at the boundary. The Kleene argument therefore does not
apply to it; termination is still guaranteed by the engine's iteration
cap (the last iterate is returned), and in practice the operator is
used on the acyclic learning path of Section~\ref{sec:learning}, where
scoring is a single fold rather than a fixpoint iteration.
Section~\ref{sec:discussion} returns to monotonicity in evidence as a
rank-robustness property in its own right.

Beyond the sketch, the properties --- monotone trajectory, boundedness,
$\varepsilon$-fixpoint reached, and exact agreement between the engine
and a reference implementation of $G$ --- are machine-checked on
1{,}050 random rule systems including seeded cycles, across the three
monotone aggregators. Property testing also exposed a real defect: a
$0.95/0.95$ cycle with a $0.05$ seed converges at a geometric rate close
to $1$ and exhausted the engine's previous iteration cap; the cap is now
sized for such cases and the example is pinned to its closed-form
fixpoint as a regression test.

\subsection{Symbolic layer and gating}

A separate crisp layer holds definite-clause facts and rules. Entailment
is decided by agenda-based forward chaining --- sound and complete for
definite clauses, linear in KB size --- differential-tested against a
truth-table oracle on random KBs; a 200-symbol chain answers in well
under a second. A hybrid engine gates confidences with entailment: a
\texttt{hard} gate zeroes non-entailed conclusions, a \texttt{soft} gate
multiplies them by a penalty, and a \texttt{constraint} mode intended to
zero out \emph{contradicted} conclusions is a documented no-op until the
symbolic language gains negation.

\subsection{What the numbers are not}\label{sec:notprob}

\pla{} has no distribution semantics: no possible-world measure, no
conditioning on evidence, and the aggregation operator \emph{chooses} an
independence stance rather than deriving one. The repository carries an
executable counterexample: two rules with confidence $0.5$ driven by the
same antecedent combine by noisy-OR to $0.75$, as if their supports were
independent, whereas a possible-world semantics would let the rules'
dependence determine the answer. Users who need that guarantee should
use ProbLog~\cite{deraedt2007}; \pla{} trades it for hand-traceable
inference.
